\chapter{Felhasználói dokumentáció}
\label{ch:user}

\section{Követelmények}

\begin{table}[H]
	\centering
	\begin{tabular}{ | m{0.25\textwidth} | m{0.65\textwidth} | }
		\hline
		\textbf{} & \textbf{Követelmény} \\
		\hline \hline
		\emph{Eszköz} & Személyi számítógép vagy laptop \\
		\hline
		\emph{Hálozati elérhetőség} &  Folyamatos stabil internetkapcsolat \\
		\hline
		\emph{Kijelző méret} & Minimum 13" \\
		\hline
		\emph{Böngésző} & Modern webböngésző (pl. Chrome vagy Safari) \\
		\hline

	\end{tabular}
	\caption{A program hardver és szoftver követelményei}
	\label{tab:example-1}
\end{table}

\section{Rövid weboldalleírás}

A webalkalmazás kezdőképernyője letisztult navigációt biztosít, ahol a játékszabályok egy dedikált aloldalon részletesen áttanulmányozhatók. A rendszer többféle hozzáférési szintet kínál: a felhasználók választhatják a regisztrációhoz kötött bejelentkezést, de lehetőség van bejelentkezés nélküli, vendégként történő játékra is. A bejelentkezési felületről a még fiókkal nem rendelkező látogatók könnyedén átnavigálhatnak a regisztrációs űrlapra. A sikeresen bejelentkezett felhasználók számára a főoldalon elérhetővé válik egy extra funkció, amely a személyes statisztikák megtekintésére szolgál.

\subsection {A játékmenet az alábbi folyamat mentén zajlik:}

\begin{itemize}
    \item \textbf{Meccskeresés:} A mérkőzés indításához vendégként a név megadása, azonosított profil esetén pedig a bejelentkezés szükséges. A párkeresési folyamat a \textbf{"Játék Keresése"} gomb aktiválásával indul, amint a rendszer megtalálja az ellenfelet, a mérkőzés automatikusan kezdetét veszi.
    \item \textbf{Játékmechanika:} A grafikus felület vizuálisan kiemeli az éppen aktív (soron lévő) játékost, a lépésjog pedig minden kártyakijátszás követően automatikusan átkerül az ellenfélhez. A kártyák kijátszása intuitív módon, egy egyszerű kattintással történik, melynek hatására a kiválasztott lap a játéktér közepén jelenik meg.
    \item \textbf{Pontozás és győzelem:} A felhasználói felület valós időben nyomon követi és megjeleníti az aktuális kártyák pontértékeit, valamint az összesített meccspontokat.
    \item \textbf{Lezárás:} A játék célja a 7 pont elérése. Amint valamelyik fél megszerzi a győzelemhez szükséges pontmennyiséget, a szoftver ezt egyértelműen jelzi a résztvevőknek, majd mindkét játékost automatikusan visszairányítja a főmenübe.
\end{itemize}

\section{Oldaltérkép és navigáció}
\label{sec:minimap_navigation}

A webes Snapszer alkalmazás felépítése a letisztult és gyors használhatóság elvét követi. Annak érdekében, hogy a játékosok a lehető leggyorsabban játékba kerüljön, a felületek közötti navigáció egy logikus, jól áttekinthető folyamatra épül. 

A \ref{fig:minimap}. ábra az alkalmazás főbb nézeteit és az azok közötti átmeneteket (felhasználói útvonalat) szemlélteti.

\subsection{Oldaltértkép}

\begin{figure}[H]
	\centering
	\includegraphics[width=0.7\textwidth]{oldal_terkep}
	\caption{Oldaltérkép}
	\label{fig:minimap}
\end{figure}

Az Oldaltérkép nevű diagram elkészítéséhez a draw.io \cite{drawio} nevű alkalmazást használtam.

\subsection{Navigáció}

A játékélmény szempontjából a navigáció az alábbi főbb állomásokból áll össze:

\begin{itemize}
    \item \textbf{Főoldal:} Az alkalmazás központi csomópontja.
    \item \textbf{Szabály oldal:} A játékszabályokat bemutató információs felület.
    \item \textbf{Bejelentkezés:} A felhasználói hitelesítés helyszíne.
    \item \textbf{Regisztráció:} Új fiók létrehozására szolgáló űrlap.
    \item \textbf{Játék oldal:} A tényleges játékmenet asztala.
    \item \textbf{Statisztika:} A bejelentkezett felhasználók személyes eredményeinek összesítője.
\end{itemize}

A következőkben az egyes oldalak részletes szerepét és felépítését mutatom be.

\subsubsection{Főoldal}
A Főoldal az alkalmazás központi csomópontja, ahonnan a rendszer összes többi funkciója közvetlenül vagy közvetve elérhető. Itt található a játék indításához és az ellenfélkereséshez szükséges felület, valamint innen navigálhatunk a szabályzat, illetve a bejelentkezés irányába is. Sikeres azonosítás után a Főoldalról válik elérhetővé a személyes statisztikák megtekintése is.

\begin{figure}[H]
	\centering
	\includegraphics[width=0.85\textwidth]{fooldal}
	\caption{Főoldal}
	\label{fig:homepage}
\end{figure}

\subsubsection{Szabály oldal}
Ezen az aloldalon a Snapszer kártyajáték részletes szabályzata olvasható. A Főoldalról közvetlenül elérhető felület célja, hogy a kezdő játékosok elsajátíthassák az alapokat, a tapasztaltabbak pedig felfrissíthessék tudásukat a lapok értékeiről és a speciális bemondásokról (például a húsz és negyven).

\begin{figure}[H]
	\centering
	\includegraphics[width=0.85\textwidth]{szabaly}
	\caption{Szabály oldal}
	\label{fig:rulepage}
\end{figure}

\subsubsection{Bejelentkezés}
A Bejelentkezés oldal szolgál a már regisztrált felhasználók hitelesítésére és profiljuk elérésére. A Főoldalról nyíló nézeten a megfelelő azonosítók (felhasználónév és jelszó) megadásával a játékos beléphet a fiókjába. Ez a felület funkcionál egyben belépési pontként a regisztrációs oldal felé is, a még fiókkal nem rendelkező látogatók számára.

\begin{figure}[H]
	\centering
	\includegraphics[width=0.85\textwidth]{bejelentkezes}
	\caption{Bejelentkezés oldal}
	\label{fig:loginpage}
\end{figure}

\subsubsection{Regisztráció}
Ahogy az oldaltérképen is látható, a Regisztráció oldal szorosan kapcsolódik a Bejelentkezéshez, és kizárólag onnan érhető el. Ezen a felületen egy letisztult űrlap kitöltésével a felhasználók létrehozhatják saját fiókjukat, amely lehetővé teszi a későbbi játékeredmények és statisztikák elmentését.

\begin{figure}[H]
	\centering
	\includegraphics[width=0.85\textwidth]{regisztracio}
	\caption{Regisztrációs oldal}
	\label{fig:registerpage}
\end{figure}

\subsubsection{Játék oldal}
A Játék oldal a mérkőzések tényleges helyszíne, maga a virtuális kártyaasztal. Ide a játékosok a Főoldalról induló sikeres meccskeresési folyamat után, emberi beavatkozás nélkül, automatikusan kerülnek át. A felületen történik a kártyák interaktív kijátszása, az aktuális pontok valós idejű követése, valamint a győzelmi feltételek teljesülése esetén a játszma lezárása.

\begin{figure}[H]
	\centering
	\includegraphics[width=0.85\textwidth]{jatek}
	\caption{Játék oldal}
	\label{fig:gamepage}
\end{figure}

\subsubsection{Statisztika}
A Statisztika egy extra funkciót kínáló felület, amely kizárólag a sikeresen bejelentkezett felhasználók számára válik elérhetővé a Főoldalról. Ezen a nézeten a játékosok nyomon követhetik a korábbi mérkőzéseik eredményeit, a lejátszott játszmák számát, valamint a győzelmek és vereségek arányát, ezáltal folyamatos visszajelzést kapva a saját teljesítményükről.

\begin{figure}[H]
	\centering
	\includegraphics[width=0.85\textwidth]{statisztika}
	\caption{Statisztika oldal}
	\label{fig:statpage}
\end{figure}

\section{A játék indítása}

A Snapszer webes alkalmazás úgy lett kialakítva, hogy a játékosok a lehető leggyorsabban, bonyolult beállítások nélkül asztalhoz ülhessenek. A játék indítása egy automatikus ellenfélkereső rendszeren keresztül történik.

\subsection{Név megadása}

A felület megnyitása után első lépésként egy tetszőleges becenév (felhasználónév) megadása szükséges, amely a játék során azonosít téged. Regisztráció vagy fiók létrehozása nem kötelező.

\begin{figure}[H]
	\centering
	\includegraphics[width=0.5\textwidth]{nev_megadasa}
	\caption{Név megadása}
	\label{fig:insert_name}
\end{figure}

\subsection{Játék keresése}

A főmenübe jutva a játék indításához csupán a "Játék indítása" (vagy Keresés) gombra kell kattintanod, miután megadtad a kívánt nevedet.

\begin{figure}[H]
	\centering
	\includegraphics[width=0.5\textwidth]{jatek_kereses}
	\caption{Játék keresés elinditása}
	\label{fig:search}
\end{figure}

\subsection{Várakozás és kezdés}

Miután megadtunk a nevünket és rákattintottunk a \textbf{"Játék Keresése"} gombra, elindul az ellenfélkeresés, és a képernyőn megjelenik egy visszaszámláló. A rendszer ekkor a háttérben egy másik, szintén várakozó játékossal próbál összepárosítani. Amint a rendszer talál egy szabad ellenfelet, a keresési képernyő eltűnik, és a játék azonnal, automatikusan kezdetét veszi az első kártyák kiosztásával.

\begin{figure}[H]
	\centering
	\includegraphics[width=0.5\textwidth]{jatek_kereses_indult}
	\caption{Játék keresés}
	\label{fig:searching}
\end{figure}

Játék keresés közben meglehet állítani a keresést az "X" gombbal.

\subsection{Játékkeresés bejelentkezett felhasználóként}
Amennyiben a játékos rendelkezik regisztrált felhasználói fiókkal és be van jelentkezve, a játék indításának folyamata még gördülékenyebb. 

Ebben az esetben a rendszer automatikusan a profilhoz tartozó nevet használja, így a név megadása lépés teljesen kimarad. A regisztrált nevünk az eddigi megszokott szövegdobozban jelenik meg. A főmenüben elegendő a \textbf{"Játék Keresése"} gombra kattintani, és azonnal kezdetét veszi az ellenfélkeresés.

Fontos különbség a vendég módhoz képest, hogy a bejelentkezett állapotban lejátszott mérkőzések eredményeit (győzelmek, vereségek) a rendszer a háttérben rögzíti. Ezek az adatok elmentésre kerülnek a felhasználó személyes statisztikájába, és lehetőséget biztosítanak a globális ranglistán való szereplésre.

\begin{figure}[H]
	\centering
	\includegraphics[width=0.5\textwidth]{bejelentkezett_kereses}
	\caption{Bejelentkezett keresés}
	\label{fig:searching_logged_in}
\end{figure}

\section{}

